\chapter{哈希、窗口、掩码：动手模拟一次``叮铃''}
\label{ch:N-hash}

\section{目标}

不写程序，用手感受：窗口滑动时摘要怎么变，何时铃响。

\section{玩具规则（极度简化）}

\begin{toybox}[只有 3 格的窗口]
用``各位数字相加再取后两位''当摘要（真 Gear 复杂得多，但感觉类似）。\\
掩码：后两位能被 5 整除 $\Rightarrow$ 叮。
\end{toybox}

\begin{figure}[htbp]
\centering
\begin{tikzpicture}[font=\small]
  \foreach \i/\v in {0/2,1/3,2/4,3/1,4/6,5/2} {
    \node[bytecell] (b\i) at (\i*1.1,0) {\v};
  }
  \draw[thick, Gen3Color] (-0.4,-0.5) rectangle (2.7,0.5);
  \node[Gen3Color, above] at (1.15,0.55) {第一步窗: 2+3+4=9};
  \node[font=\footnotesize, below] at (1.15,-0.85) {9 不能被5整除 $\Rightarrow$ 不叮};
\end{tikzpicture}
\caption{第一步：不切。}
\label{fig:N-toy1}
\end{figure}

\begin{figure}[htbp]
\centering
\begin{tikzpicture}[font=\small]
  \foreach \i/\v in {0/2,1/3,2/4,3/1,4/6,5/2} {
    \node[bytecell] (b\i) at (\i*1.1,0) {\v};
  }
  \draw[thick, Gen3Color] (0.7,-0.5) rectangle (3.8,0.5);
  \node[Gen3Color, above] at (2.25,0.55) {滑一步: 3+4+1=8};
  \node[font=\footnotesize, below] at (2.25,-0.85) {仍不叮};
\end{tikzpicture}
\caption{第二步：仍不切。}
\label{fig:N-toy2}
\end{figure}

\begin{figure}[htbp]
\centering
\begin{tikzpicture}[font=\small]
  \foreach \i/\v in {0/2,1/3,2/4,3/1,4/6,5/2} {
    \node[bytecell] (b\i) at (\i*1.1,0) {\v};
  }
  \draw[thick, OkGreen] (1.8,-0.5) rectangle (4.9,0.5);
  \node[OkGreen, above] at (3.35,0.55) {再滑: 4+1+6=11 $\to$ 后两位像 11};
  \draw[cutmark] (4.9,0.7)--(4.9,-0.7);
  \node[font=\footnotesize, CutRed, below] at (4.9,-0.95) {若规则说叮：这里成为候选刀口};
\end{tikzpicture}
\caption{真算法铃响还要过第二关；这里只演示``叮''的感觉。}
\label{fig:N-toy3}
\end{figure}

\section{和真 Hom 的对应}

\begin{center}
\begin{tabular}{@{}ll@{}}
\toprule
玩具 & 真 Hom \\
\midrule
三位数字相加 & Gear 滚动指纹 $h$ \\
能被5整除 & $(h$ 按位与 mask$)=0$ \\
滑窗 & 固定宽度 window \\
\bottomrule
\end{tabular}
\end{center}

\begin{pitfall}[别被公式吓到]
看到 $h\leftarrow 2h+G[b_{in}]-G[b_{out}]$，只记：\\
\textbf{出一格、进一格，摘要用加减快速改}，不必每次重算整窗。
\end{pitfall}
